\documentclass[12pt]{article}
\usepackage[T1]{fontenc}
\usepackage[utf8]{inputenc}
\usepackage{lmodern}
\usepackage{textcomp}
\usepackage{enumitem}
\usepackage{geometry}
\geometry{margin=1in}
\begin{document}
\begin{center}
Answer Key - Constitutional Law - Graduate Level Assessment 24 \
\small (Answers are ordered exactly as in the question paper.)
\end{center}

\section*{Multiple Choice}
\begin{enumerate}[label=\arabic*.]
\item \textbf{Answer:} A
\\ \textit{Options:} A) The law is the only guide for ethical behaviour.; B) The law dictates what moral values should affect our ethical behaviour.; C) Morality plays no role in the concept of law.; D) Moral arguments operate only in hard cases.; E) Morality only applies to specific, complex cases.
\\ \textit{Note:} Dworkin argues that legal principles are not only rules to be followed but also serve as a comprehensive framework that guides individuals in making ethical decisions, thereby positioning law as the primary reference point for determining moral conduct.
\item \textbf{Answer:} A
\\ \textit{Options:} A) Dudgeon v UK (1981); B) Marckx v Belgium (1979); C) Osman v UK ( 1998); D) Tyrer v UK (1978); E) McCann v UK (1995)
\\ \textit{Note:} The case of Dudgeon v UK (1981) is correct because it was the first instance where the European Court of Human Rights articulated the 'margin of appreciation' doctrine, allowing states discretion in how they implement certain rights under the European Convention on Human Rights, thereby balancing individual rights with national interests.
\item \textbf{Answer:} A
\\ \textit{Options:} A) Treaties do no create obligations or rights for third States without their consent; B) Treaties create both obligations and rights for third States; C) Treaties do not create any rights for third States, even when the latter consent.; D) Treaties may create only obligations for third States; E) Treaties create obligations for third States without their consent
\\ \textit{Note:} The principle of pacta tertiis nec nocent nec prosunt establishes that treaties do not impose obligations or confer rights on third States unless those States have explicitly consented to be bound by the treaty.
\item \textbf{Answer:} C
\\ \textit{Options:} A) Related Doctrine of Merger; B) Contingent remainder; C) Doctrine of Worthier Title; D) Doctrine of Escheat; E) Vested remainder
\\ \textit{Note:} The Doctrine of Worthier Title holds that if a grantor attempts to create a future interest in their heirs, that interest is deemed invalid and instead results in a reversion back to the grantor, reflecting the principle that the grantor's intent is prioritized over the heirs.
\item \textbf{Answer:} B
\\ \textit{Options:} A) His concept of status is misrepresented.; B) It is misinterpreted as a prediction.; C) His idea is considered inapplicable to Western legal systems.; D) It is incorrectly related to economic progression.; E) It is taken literally.
\\ \textit{Note:} The correct answer highlights that the aphorism is often misinterpreted as suggesting a linear trajectory towards increased individualism in societal structures, rather than recognizing it as a descriptive observation of historical legal evolution that does not necessarily predict future outcomes.
\end{enumerate}

\section*{True / False}
\begin{enumerate}[label=\arabic*.]
\setcounter{enumi}{5}
\item \textbf{Answer:} False
\\ \textit{Options:} A) True; B) False
\\ \textit{Note:} A warranty of merchantability ensures that a product is fit for its ordinary purpose and of average quality, but it does not guarantee that it will meet all of the buyer's specific needs or preferences without further guidance.
\item \textbf{Answer:} True
\\ \textit{Options:} A) True; B) False
\\ \textit{Note:} The original position, as conceived by John Rawls, emphasizes the principles of justice and fairness that prioritize the equitable distribution of resources and opportunities rather than the mere accumulation of wealth, which can lead to societal inequalities and undermine compassion.
\item \textbf{Answer:} True
\\ \textit{Options:} A) True; B) False
\\ \textit{Note:} Passive personality jurisdiction is true because it allows a state to exercise jurisdiction over a foreign offender when the victim of the crime is a national of that state, emphasizing the importance of the victim's nationality in establishing legal authority.
\item \textbf{Answer:} False
\\ \textit{Options:} A) True; B) False
\\ \textit{Note:} Bentham argued that law is a product of human legislation rather than solely derived from moral principles or societal needs, emphasizing the distinction between law and morality.
\item \textbf{Answer:} False
\\ \textit{Options:} A) True; B) False
\\ \textit{Note:} The statement is false because solicitation is typically defined as the act of encouraging or enticing another person to commit a crime, which does not necessarily require an agreement between multiple parties, distinguishing it from conspiracy.
\end{enumerate}

\section*{Long Form Response}
\begin{enumerate}[label=\arabic*.]
\setcounter{enumi}{10}
\item \textbf{Answer:} One of the strongest arguments against ethical relativism's stance on human rights is its inability to provide a solid foundation for universal human rights, which are often seen as inherent and inalienable. Ethical relativism posits that moral truths are context- dependent, leading to the conclusion that human rights can vary significantly across cultures. This perspective undermines the cognitivist view, which asserts that moral statements can be true or false and that there are objective standards by which to evaluate them. By emphasizing cognitivism, proponents argue for a universal framework of human rights that transcends cultural differences, asserting that certain rights-such as the right to life and freedom from torture-are fundamental and should be recognized regardless of cultural context. Thus, the relativist position is critiqued for failing to accommodate the notion of shared human dignity and the moral obligations that arise from it.
\\ \textit{Note:} correct because it highlights that ethical relativism's context-dependent nature challenges the universality of human rights, which contradicts cognitivism's assertion of objective moral truths, thus undermining the foundation for claiming inherent human rights.
\item \textbf{Answer:} In auction settings, the legal implications of bid retraction are significant, particularly illustrated by the context of Cross \& Black's liquidation. The principle that a bidder may retract their bid at any time until the falling of the hammer underscores the fluidity and conditional nature of bids within auction law. This provision protects bidders' rights by allowing them to reassess their commitments and avoid unforeseen liabilities, fostering a competitive bidding environment. However, it also raises concerns about the integrity of the bidding process, as frequent retractions can undermine trust and lead to instability in auction outcomes. Ultimately, while the right to retract bids serves as a safeguard for bidders, it necessitates a careful balance to maintain the legitimacy of the auction framework.
\\ \textit{Note:} correct because it emphasizes the legal principle that allows bidders to retract their bids until the auction concludes, thereby protecting their rights and enabling a competitive bidding atmosphere, as illustrated by the context of Cross \& Black's liquidation.
\end{enumerate}

\vfill
\noindent\textit{End of Answer Key}
\end{document}